% GENERATED FILE. Edit canonical lessons, then run npm run generate:books.
% canonical-source-hash: fnv1a64:3dc312b6b9bb4eca
% canonical-lessons: JA-C479-nakayoshi, JA-C479-nakanaori-suru, JA-C479-shitashii, JA-C479-mimamoru, JA-C479-uyamau, JA-R479-review-474-479, JA-R479-review-final-479

\chapter{Close Ties}
\label{ch:ja-close-ties}
\hlchaptermodality{479}

\begin{chapteropening}
\textbf{By the end of this chapter:} \emph{I can talk about close friends and family ties.}

Complete the last lesson of chapter 479: I can talk about close friends and family ties.
\end{chapteropening}

\section[\texorpdfstring{\ja{なかよし} (nakayoshi)}{なかよし (nakayoshi)}]{\ja{なかよし} (nakayoshi) — close friends; being close}

\label{lesson:JA-C479-nakayoshi}



\begin{quote}
Before the new one: say the Japanese for to tease, then the Japanese for to imitate.
\end{quote}

\subsection*{You'll want to know: \ja{なかよし}}
\textbf{\ja{なかよし}} — \emph{nakayoshi} — \textquotedblleft{}close friends; being close\textquotedblright{}.

\begin{quote}
\textbf{\ja{いもうとと} \ja{となりの} \ja{こどもは} \ja{なかよしです。}} — \emph{imōto to tonari no kodomo wa nakayoshi desu} — \textquotedblleft{}My younger sister and the child next door are close friends.\textquotedblright{}
\end{quote}

\begin{cousinweb}
\textbf{\ja{なかよし}} (\emph{nakayoshi}) is a native Japanese compound (wago), inherited rather than borrowed from Chinese: \emph{naka}, \textquotedblleft{}relationship\textquotedblright{} and \emph{yoshi}, \textquotedblleft{}good\textquotedblright{} (JMdict, KANJIDIC2).
\end{cousinweb}

\begin{grammarlens}[title={What you've built}]
One more word for describing what you lived and felt.
\end{grammarlens}

\subsection*{Your turn}
\begin{itemize}
\raggedright
  \item \emph{Say it:} \emph{nakayoshi}
  \item \emph{Say it:} \emph{nakayoshi}, once more
  \item \emph{Say it:} \emph{imōto to tonari no kodomo wa nakayoshi desu} — My younger sister and the child next door are close friends.
  \item \emph{Recall:} say the Japanese for to tease, then the Japanese for to imitate, then say \emph{nakayoshi} again
\end{itemize}

\begin{tcolorbox}[breakable,colback=teal!4,colframe=teal!35!black,title={Before you move on}]
What does \ja{なかよし} mean? (\textquotedblleft{}Close friends; being close\textquotedblright{}.) Say it once more.
\end{tcolorbox}
\section[\texorpdfstring{\ja{なかなおりする} (nakanaori suru)}{なかなおりする (nakanaori suru)}]{\ja{なかなおりする} (nakanaori suru) — to make up (after a quarrel)}

\label{lesson:JA-C479-nakanaori-suru}



\begin{quote}
Before the new one: say the Japanese for to imitate, then the Japanese for close friends.
\end{quote}

\subsection*{You'll want to know: \ja{なかなおりする}}
\textbf{\ja{なかなおりする}} — \emph{nakanaori suru} — \textquotedblleft{}to make up (after a quarrel)\textquotedblright{}.

\begin{quote}
\textbf{\ja{けんかの} \ja{あとで、なかなおりしました。}} — \emph{kenka no ato de, nakanaori shimashita} — \textquotedblleft{}After the quarrel we made up.\textquotedblright{}
\end{quote}

Polite \textbf{\ja{なかなおりします}}, \emph{nakanaorishimasu}; a \ja{する} verb. Past (polite): \textbf{\ja{ふたりが} \ja{なかなおりしました。}}, \emph{futari ga nakanaori shimashita}.

\begin{cousinweb}
\textbf{\ja{なかなおりする}} (\emph{nakanaori suru}) is a native Japanese compound (wago), inherited rather than borrowed from Chinese: \emph{naka}, \textquotedblleft{}relationship\textquotedblright{} and \emph{naori}, \textquotedblleft{}mending (naoru)\textquotedblright{} (JMdict, KANJIDIC2).
\end{cousinweb}

\begin{grammarlens}[title={What you've built}]
One more word for describing what you lived and felt.
\end{grammarlens}

\subsection*{Your turn}
\begin{itemize}
\raggedright
  \item \emph{Say it:} \emph{nakanaori suru}
  \item \emph{Say it:} \emph{nakanaori suru}, once more
  \item \emph{Say it:} \emph{kenka no ato de, nakanaori shimashita} — After the quarrel we made up.
  \item \emph{Recall:} say the Japanese for to imitate, then the Japanese for close friends, then say \emph{nakanaori suru} again
\end{itemize}

\begin{tcolorbox}[breakable,colback=teal!4,colframe=teal!35!black,title={Before you move on}]
What does \ja{なかなおりする} mean? (\textquotedblleft{}To make up (after a quarrel)\textquotedblright{}.) Say it once more.
\end{tcolorbox}
\section[\texorpdfstring{\ja{したしい} (shitashii)}{したしい (shitashii)}]{\ja{したしい} (shitashii) — close to someone; friendly}

\label{lesson:JA-C479-shitashii}



\begin{quote}
Before the new one: say the Japanese for close friends, then the Japanese for to make up.
\end{quote}

\subsection*{You'll want to know: \ja{したしい}}
\textbf{\ja{したしい}} — \emph{shitashii} — \textquotedblleft{}close to someone; friendly\textquotedblright{}.

\begin{quote}
\textbf{\ja{あには} \ja{どうりょうと} \ja{したしいです。}} — \emph{ani wa dōryō to shitashii desu} — \textquotedblleft{}My elder brother is close to his colleagues.\textquotedblright{}
\end{quote}

An i-adjective.

\begin{cousinweb}
\textbf{\ja{したしい}} (\emph{shitashii}) is a native Japanese word (wago), inherited rather than borrowed from Chinese; its kanji spelling was attached to it later (JMdict).
\end{cousinweb}

\begin{grammarlens}[title={What you've built}]
One more word for describing what you lived and felt.
\end{grammarlens}

\subsection*{Your turn}
\begin{itemize}
\raggedright
  \item \emph{Say it:} \emph{shitashii}
  \item \emph{Say it:} \emph{shitashii}, once more
  \item \emph{Say it:} \emph{ani wa dōryō to shitashii desu} — My elder brother is close to his colleagues.
  \item \emph{Recall:} say the Japanese for close friends, then the Japanese for to make up, then say \emph{shitashii} again
\end{itemize}

\begin{tcolorbox}[breakable,colback=teal!4,colframe=teal!35!black,title={Before you move on}]
What does \ja{したしい} mean? (\textquotedblleft{}Close to someone; friendly\textquotedblright{}.) Say it once more.
\end{tcolorbox}
\section[\texorpdfstring{\ja{みまもる} (mimamoru)}{みまもる (mimamoru)}]{\ja{みまもる} (mimamoru) — to watch over}

\label{lesson:JA-C479-mimamoru}



\begin{quote}
Before the new one: say the Japanese for to make up, then the Japanese for close to someone.
\end{quote}

\subsection*{You'll want to know: \ja{みまもる}}
\textbf{\ja{みまもる}} — \emph{mimamoru} — \textquotedblleft{}to watch over\textquotedblright{}.

\begin{quote}
\textbf{\ja{りょうしんは} \ja{こどもを} \ja{みまもっています。}} — \emph{ryōshin wa kodomo o mimamotte imasu} — \textquotedblleft{}The parents are watching over the child.\textquotedblright{}
\end{quote}

Polite \textbf{\ja{みまもります}}, \emph{mimamorimasu}; a godan verb. Past (polite): \textbf{\ja{ははが} \ja{こどもを} \ja{みまもりました。}}, \emph{haha ga kodomo o mimamorimashita}.

\begin{cousinweb}
\textbf{\ja{みまもる}} (\emph{mimamoru}) is a native Japanese compound (wago), inherited rather than borrowed from Chinese: \emph{mi}, \textquotedblleft{}seeing (miru)\textquotedblright{} and \emph{mamoru}, \textquotedblleft{}to protect\textquotedblright{} (JMdict, KANJIDIC2).
\end{cousinweb}

\begin{grammarlens}[title={What you've built}]
One more word for describing what you lived and felt.
\end{grammarlens}

\subsection*{Your turn}
\begin{itemize}
\raggedright
  \item \emph{Say it:} \emph{mimamoru}
  \item \emph{Say it:} \emph{mimamoru}, once more
  \item \emph{Say it:} \emph{ryōshin wa kodomo o mimamotte imasu} — The parents are watching over the child.
  \item \emph{Recall:} say the Japanese for to make up, then the Japanese for close to someone, then say \emph{mimamoru} again
\end{itemize}

\begin{tcolorbox}[breakable,colback=teal!4,colframe=teal!35!black,title={Before you move on}]
What does \ja{みまもる} mean? (\textquotedblleft{}To watch over\textquotedblright{}.) Say it once more.
\end{tcolorbox}
\section[\texorpdfstring{\ja{うやまう} (uyamau)}{うやまう (uyamau)}]{\ja{うやまう} (uyamau) — to revere; to honour}

\label{lesson:JA-C479-uyamau}



\begin{quote}
Before the new one: say the Japanese for close to someone, then the Japanese for to watch over.
\end{quote}

\subsection*{You'll want to know: \ja{うやまう}}
\textbf{\ja{うやまう}} — \emph{uyamau} — \textquotedblleft{}to revere; to honour\textquotedblright{}.

\begin{quote}
\textbf{\ja{みんな} \ja{あの} \ja{せんせいを} \ja{うやまっています。}} — \emph{minna ano sensei o uyamatte imasu} — \textquotedblleft{}Everyone respects that teacher.\textquotedblright{}
\end{quote}

Polite \textbf{\ja{うやまいます}}, \emph{uyamaimasu}; a godan verb. Past (polite): \textbf{\ja{がくせいが} \ja{せんせいを} \ja{うやまいました。}}, \emph{gakusei ga sensei o uyamaimashita}.

\begin{cousinweb}
\textbf{\ja{うやまう}} (\emph{uyamau}) is a native Japanese word (wago), inherited rather than borrowed from Chinese; its kanji spelling was attached to it later (JMdict).
\end{cousinweb}

\begin{grammarlens}[title={What you've built}]
One more word for describing what you lived and felt.
\end{grammarlens}

\subsection*{Your turn}
\begin{itemize}
\raggedright
  \item \emph{Say it:} \emph{uyamau}
  \item \emph{Say it:} \emph{uyamau}, once more
  \item \emph{Say it:} \emph{minna ano sensei o uyamatte imasu} — Everyone respects that teacher.
  \item \emph{Recall:} say the Japanese for close to someone, then the Japanese for to watch over, then say \emph{uyamau} again
\end{itemize}

\begin{tcolorbox}[breakable,colback=teal!4,colframe=teal!35!black,title={Before you move on}]
What does \ja{うやまう} mean? (\textquotedblleft{}To revere; to honour\textquotedblright{}.) Say it once more.
\end{tcolorbox}
\section[(dialogue)]{Looking back: from Getting the Food Ready to Close Ties}

\label{lesson:JA-R479-review-474-479}



\begin{quote}
Say the Japanese for to squeeze and to mash.
\end{quote}

\subsection*{Your turn}
Say these aloud:

\begin{itemize}
\raggedright
  \item to squeeze, to mash, to add, to tear into pieces, to dilute
  \item a fridge, tableware, a tablespoon, a small plate, flavouring
  \item a cousin, twins, an only child, the eldest son, the eldest daughter
  \item a parent, parent and child, the youngest child, a bother, upbringing
  \item to pester for, to spoil a child, to bully, to tease, to imitate
  \item close friends, to make up, close to someone, to watch over, to revere
\end{itemize}

\begin{tcolorbox}[breakable,colback=teal!4,colframe=teal!35!black,title={Before you move on}]
To squeeze? (\textbf{\ja{しぼる}}, \emph{shiboru}.) To mash? (\textbf{\ja{つぶす}}, \emph{tsubusu}.) To add? (\textbf{\ja{くわえる}}, \emph{kuwaeru}.) To tear into pieces? (\textbf{\ja{ちぎる}}, \emph{chigiru}.) To dilute? (\textbf{\ja{うすめる}}, \emph{usumeru}.) A fridge? (\textbf{\ja{れいぞうこ}}, \emph{reizōko}.) Tableware? (\textbf{\ja{しょっき}}, \emph{shokki}.) A tablespoon? (\textbf{\ja{おおさじ}}, \emph{ōsaji}.) A small plate? (\textbf{\ja{こざら}}, \emph{kozara}.) Flavouring? (\textbf{\ja{あじつけ}}, \emph{ajitsuke}.) A cousin? (\textbf{\ja{いとこ}}, \emph{itoko}.) Twins? (\textbf{\ja{ふたご}}, \emph{futago}.) An only child? (\textbf{\ja{ひとりっこ}}, \emph{hitorikko}.) The eldest son? (\textbf{\ja{ちょうなん}}, \emph{chōnan}.) The eldest daughter? (\textbf{\ja{ちょうじょ}}, \emph{chōjo}.) A parent? (\textbf{\ja{おや}}, \emph{oya}.) Parent and child? (\textbf{\ja{おやこ}}, \emph{oyako}.) The youngest child? (\textbf{\ja{すえっこ}}, \emph{suekko}.) A bother? (\textbf{\ja{めんどう}}, \emph{mendō}.) Upbringing? (\textbf{\ja{しつけ}}, \emph{shitsuke}.) To pester for? (\textbf{\ja{ねだる}}, \emph{nedaru}.) To spoil a child? (\textbf{\ja{あまやかす}}, \emph{amayakasu}.) To bully? (\textbf{\ja{いじめる}}, \emph{ijimeru}.) To tease? (\textbf{\ja{からかう}}, \emph{karakau}.) To imitate? (\textbf{\ja{まねる}}, \emph{maneru}.) Close friends? (\textbf{\ja{なかよし}}, \emph{nakayoshi}.) To make up? (\textbf{\ja{なかなおりする}}, \emph{nakanaori suru}.) Close to someone? (\textbf{\ja{したしい}}, \emph{shitashii}.) To watch over? (\textbf{\ja{みまもる}}, \emph{mimamoru}.) To revere? (\textbf{\ja{うやまう}}, \emph{uyamau}.)
\end{tcolorbox}
\section[(dialogue)]{Close Ties — once more}

\label{lesson:JA-R479-review-final-479}



\begin{quote}
Say the Japanese for close friends and to make up.
\end{quote}

\subsection*{Your turn}
\begin{itemize}
\raggedright
  \item \emph{Say it:} close friends, to make up, close to someone, to watch over, to revere
\end{itemize}

\begin{tcolorbox}[breakable,colback=teal!4,colframe=teal!35!black,title={Before you move on}]
Close friends? (\textbf{\ja{なかよし}}, \emph{nakayoshi}.) To make up? (\textbf{\ja{なかなおりする}}, \emph{nakanaori suru}.) Close to someone? (\textbf{\ja{したしい}}, \emph{shitashii}.) To watch over? (\textbf{\ja{みまもる}}, \emph{mimamoru}.) To revere? (\textbf{\ja{うやまう}}, \emph{uyamau}.)
\end{tcolorbox}
